\section{Results}
\label{sec:results}

\subsection{Direction Quality Validation}
\label{sec:direction_quality}

Before testing composition, we validate that each concept direction captures
meaningful information. \Tabref{tab:probing} reports the best-layer linear
probing accuracy for all 10 concepts.

\begin{table}[t]
    \centering
    \caption{Linear probing accuracy for individual concept directions. All
    concepts except true/false achieve $\geq$72.5\% accuracy, validating that
    extracted directions encode meaningful information. Concepts peak at
    different layers, ranging from layer~3 (formal/casual) to layer~31
    (fruit/vegetable).}
    \begin{tabular}{@{}lcc@{}}
        \toprule
        \textbf{Concept} & \textbf{Best Acc.} & \textbf{Best Layer} \\
        \midrule
        formal/casual     & \textbf{1.000} & 3  \\
        singular/plural   & 0.950 & 24 \\
        happy/sad         & 0.950 & 17 \\
        big/small         & 0.925 & 14 \\
        sentiment (+/$-$) & 0.850 & 11 \\
        past/present      & 0.825 & 11 \\
        fruit/vegetable   & 0.800 & 31 \\
        hot/cold          & 0.775 & 21 \\
        mammal/bird       & 0.725 & 20 \\
        true/false        & 0.650 & 15 \\
        \bottomrule
    \end{tabular}
    \label{tab:probing}
\end{table}

All concepts except true/false ($0.650$) exceed 72.5\% accuracy, confirming
that the mean-difference extraction method produces meaningful directions. The
wide range in optimal layers---from layer~3 for formal/casual to layer~31 for
fruit/vegetable---indicates that different concepts are resolved at different
depths, consistent with prior observations that syntactic features emerge
earlier than semantic ones~\citep{stolfo2024improving}.

\subsection{Pairwise Orthogonality}
\label{sec:orthogonality}

\Tabref{tab:orthogonality} reports cosine similarities between all 10 tested
direction pairs at layer~20.

\begin{table}[t]
    \centering
    \caption{Pairwise cosine similarity between concept directions at layer~20.
    Most pairs are far from orthogonal: the median $\abscossim$ is 0.91. Only 3
    of 10 pairs have $\abscossim < 0.3$. Anti-parallel pairs (negative cosine)
    are as problematic for composition as parallel pairs.}
    \begin{tabular}{@{}llcc@{}}
        \toprule
        \textbf{Pair} & \textbf{Category} & $\cossim$ & $\abscossim$ \\
        \midrule
        true/false + sentiment        & entangled     & $-0.994$ & 0.994 \\
        past/pres.\ + sing./plur.\    & syntactic     & $-0.954$ & 0.954 \\
        big/small + sentiment         & semantic      & $+0.937$ & 0.937 \\
        big/small + hot/cold          & semantic      & $-0.923$ & 0.923 \\
        past/pres.\ + big/small       & cross-domain  & $+0.919$ & 0.919 \\
        true/false + formal/casual    & cross-domain  & $+0.897$ & 0.897 \\
        sentiment + formal/casual     & independent   & $-0.888$ & 0.888 \\
        \midrule
        mammal/bird + sing./plur.\    & hierarchical  & $+0.185$ & 0.185 \\
        sentiment + happy/sad         & entangled     & $-0.153$ & 0.153 \\
        mammal/bird + fruit/veg.\     & hierarchical  & $-0.106$ & \textbf{0.106} \\
        \bottomrule
    \end{tabular}
    \label{tab:orthogonality}
\end{table}

{\bf The central surprise is the degree of non-orthogonality.} Seven of 10
pairs have $\abscossim > 0.88$, and the median is $0.908$. Only 3 pairs
(mammal/bird + fruit/vegetable, sentiment + happy/sad, mammal/bird +
singular/plural) are near-orthogonal ($\abscossim < 0.3$). Notably, seemingly
unrelated concepts like past/present tense and big/small show high alignment
($\cossim = 0.919$), suggesting that the mean-difference extraction method
captures substantial shared variance beyond the target concept.

\subsection{Probing-Based Composition}
\label{sec:probing_results}

Information preservation---the ratio of composed-direction probing accuracy to
individual-direction accuracy---is uniformly high across all categories
(\tabref{tab:probing_composition}).

\begin{table}[t]
    \centering
    \caption{Probing-based composition preservation by expected composability
    category. Preservation is uniformly high (${\sim}0.95$) and does not
    differ significantly across categories (ANOVA $F = 0.86$, $p = 0.46$).
    Counter-intuitively, ``low'' expected composability pairs show the
    highest preservation.}
    \begin{tabular}{@{}lcc@{}}
        \toprule
        \textbf{Expected} & \textbf{Mean Preservation} & \textbf{95\% CI} \\
        \midrule
        High   & 0.926 & [0.882, 0.966] \\
        Medium & 0.958 & [0.911, 0.995] \\
        Low    & 0.983 & [0.966, 1.000] \\
        \midrule
        Overall & \textbf{0.954} & [0.925, 0.980] \\
        \bottomrule
    \end{tabular}
    \label{tab:probing_composition}
\end{table}

One-way ANOVA yields $F = 0.86$, $p = 0.46$---the expected composability
categories do not predict probing-based composition. The counter-intuitive
result that ``low'' expected pairs show the \emph{highest} preservation arises
because the 1D projection metric is permissive: even anti-parallel directions
sum to a vector that partially captures each concept's variance. This metric
alone would suggest that composition works universally, which our steering
results contradict.

\subsection{Steering-Based Composition}
\label{sec:steering_results}

Steering reveals much sharper differentiation. \Tabref{tab:steering} reports
component preservation and shift alignment for all 10 pairs.

\begin{table}[t]
    \centering
    \caption{Steering-based composition results at layer~20. Component
    preservation measures how much of each individual direction's logit
    shift is retained in the composed vector. Shift alignment is the cosine
    similarity between the composed shift and the sum of individual shifts.
    Bold values indicate component preservation $\geq 0.65$.}
    \resizebox{\textwidth}{!}{%
    \begin{tabular}{@{}llccccc@{}}
        \toprule
        \textbf{Pair} & \textbf{Expected} & $\abscossim$ &
        \textbf{Comp.\ A} & \textbf{Comp.\ B} & \textbf{Shift Align.} \\
        \midrule
        mammal/bird + fruit/veg.\   & high   & 0.11 & \textbf{0.769} & \textbf{0.686} & 0.994 \\
        mammal/bird + sing./plur.\  & high   & 0.19 & \textbf{0.941} & 0.490 & 0.998 \\
        big/small + sentiment       & medium & 0.94 & \textbf{0.866} & \textbf{0.821} & 0.968 \\
        past/pres.\ + big/small     & medium & 0.92 & \textbf{0.814} & \textbf{0.879} & 0.989 \\
        true/false + formal/casual  & medium & 0.90 & 0.640          & \textbf{0.932} & 0.989 \\
        sentiment + formal/casual   & high   & 0.89 & 0.193          & \textbf{0.872} & 0.935 \\
        big/small + hot/cold        & medium & 0.92 & \textbf{0.660} & 0.374          & 0.863 \\
        past/pres.\ + sing./plur.\  & medium & 0.95 & 0.235          & 0.537          & 0.780 \\
        sentiment + happy/sad       & low    & 0.15 & 0.200          & \textbf{0.958} & 0.984 \\
        true/false + sentiment      & low    & 0.99 & 0.116          & 0.218          & 0.441 \\
        \bottomrule
    \end{tabular}
    }
    \label{tab:steering}
\end{table}

Several patterns emerge from the steering results:

\para{Near-orthogonal pairs compose best.}
The mammal/bird + fruit/vegetable pair ($\abscossim = 0.11$) preserves both
components well (0.77, 0.69) with near-perfect shift alignment (0.994). This
confirms the prediction from \citet{park2023linear} that causally separable,
hierarchically related concepts occupy orthogonal subspaces amenable to
composition.

\para{Anti-parallel pairs fail catastrophically.}
The true/false + sentiment pair ($\abscossim = 0.99$, anti-parallel) produces
the worst steering outcome: component preservation of 0.12 and 0.22, with a
shift alignment of only 0.44. The two directions effectively cancel each other,
producing a composed vector that steers in neither intended direction.

\para{Parallel alignment does not always prevent composition.}
Surprisingly, the big/small + sentiment pair ($\abscossim = 0.94$, parallel)
preserves both components well (0.87, 0.82). When directions are parallel
(same sign), their sum reinforces rather than cancels, and both logit shifts can
be partially preserved.

\para{Component asymmetry is pervasive.}
In 7 of 10 pairs, the two components differ in preservation by more than 0.2.
The direction with larger magnitude or stronger individual steering effect tends
to dominate---\eg in sentiment + happy/sad, the more specific happy/sad
direction dominates (0.96) while sentiment is nearly lost (0.20).

\subsection{Layer-Wise Stability}
\label{sec:layerwise}

We measure probing-based composition preservation across 6 layers
(\tabref{tab:layers}). Preservation is remarkably stable from layer~5 through
layer~25, with a modest decline at the final layer~31.

\begin{table}[t]
    \centering
    \caption{Mean probing-based composition preservation across layers. Values
    are mean $\pm$ standard deviation across all 10 pairs. Composition quality
    is stable across layers, with a slight decline at the final layer,
    consistent with findings from \citet{guo2025quantifying}.}
    \begin{tabular}{@{}cc@{}}
        \toprule
        \textbf{Layer} & \textbf{Mean Preservation $\pm$ SD} \\
        \midrule
        5  & $0.943 \pm 0.067$ \\
        10 & $0.956 \pm 0.081$ \\
        15 & $0.944 \pm 0.056$ \\
        20 & $0.954 \pm 0.045$ \\
        25 & $0.952 \pm 0.046$ \\
        31 & $0.906 \pm 0.103$ \\
        \bottomrule
    \end{tabular}
    \label{tab:layers}
\end{table}

The stability across layers suggests that composition properties are a
feature of the residual stream's additive structure rather than
layer-specific processing, consistent with the view of the residual
stream as an information bus~\citep{nanda2023emergent}. The decline at
layer~31 mirrors \citeauthor{guo2025quantifying}'s~(\citeyear{guo2025quantifying})
finding that compositionality drops at the top layer.
